\documentclass[../main.tex]{subfiles}
\begin{document}
\chapter{2026-2027学年微积分（B）（上）期中考试}

\section{基本计算题(每小题 6 分，共 60 分)}
\begin{enumerate}
    \item 求极限$l=\lim_{n\to\infty}\frac{1!+2!+3!+\cdots+n!}{n!}$.
          \vspace{8em}

    \item 求极限$l=\lim_{x\to2}\frac{(1-\sqrt{x-1})(1-\sqrt[3]{x-1})\cdots(1-\sqrt[n]{x-1})}{(x-2)^{n-1}}$.
          \vspace{8em}

    \item 求极限$l=\lim_{x\to0}\frac{\tan^{2}x}{1-\cos x^{2}}\left(1-\sqrt{\cos x}\right)$.
          \vspace{7em}

    \item 确定极限$\lim_{x\to\infty}\left(\frac{x^{2}+1}{x+1}-ax-b\right)=\frac{1}{2}$中的参数$a,b$.
          \vspace{7em}

    \item 当$x\to0$时，求无穷小量$\sin(\tan x)-\tan(\tan x)$的主部和阶数.
          \vspace{9em}

    \item 设$f(x)=
              \begin{cases}
                  \frac{2}{\pi}\arctan\frac{1}{x-1}, & x>1,    \\
                  \mathrm{e}^{x},                    & x\leq1,
              \end{cases}$讨论$f(x)$的可导性，可导时求$f'(x)$.
          \vspace{9em}

    \item 若$y=y(x)$由极坐标方程$r(\theta)=\mathrm{e}^{\theta}$确定，求$\frac{\mathrm{d}^{2}y}{\mathrm{d}
                  x^{2}}$.
          \vspace{9em}

    \item 求$f(x)=\frac{x^{2}-1}{x^{2}-3x+2}+\sin\frac{1}{x}$的间断点，并判断每个间断点的类型.
          \vspace{9em}

    \item 设函数$f(x)=
              \begin{cases}
                  x^{\alpha}\sin\frac{1}{x}, & x\neq0, \\
                  0,                         & x=0,
              \end{cases}$其中$\alpha$为实数.问：(1) $\alpha$为何值时，函数$f(x)$在$x=0$处连续？(2)
          $\alpha$为何值时，函数$f(x)$在$x=0$处可导？
          \vspace{9em}

    \item 设函数$y=\left[f(x^{2})\right]^{\frac{1}{x}}$，其中$f(x)$是可微的正值函数，且$f(1)=2$，$f'(1)=1$，求$\left.
              \mathrm{d}y\right|_{x=1}$.
          \vspace{9em}
\end{enumerate}

\section{综合题(每小题 6 分，共 30 分)}
\begin{enumerate}
    \setcounter{enumi}{10}
    \item 设函数$f(x)$在$x=0$处二阶可导，且
          \[
              \lim_{x\to0}\frac{\sin x+\frac{f(x)}{x}}{x}=3,
          \]
          求$f(0)$，$f'(0)$，$\lim_{x\to0}\frac{f(x)}{x^{2}}$.
          \vspace{10em}

    \item 若函数$f(x)$严格单调递增且处处二阶可导，且对于任意实数$x$，$f(x)$满足方程$f^{3}(x)+f(x)=x+2$，$f(0)=1$，
          求$y=f(x)$的反函数$x=g(y)$在$y=1$处的二阶导数$\left.\frac{\mathrm{d}^{2}x}{\mathrm{d}y^{2}}\right|_{y=1}$.
          \vspace{10em}

    \item 水自高为$24\mathrm{cm}$、底半径为$6\mathrm{cm}$的圆锥形漏斗（圆锥底部为圆面，尖朝上）流入半径为$6
              \mathrm{cm}$的圆柱形桶内.已知漏斗水深$16\mathrm{cm}$时，水面下降的速度为$2\mathrm{cm/s}$，求此时桶中水面上升的速度.
          \vspace{12em}

    \item 设$x_{1}=1$，$x_{n+1}=1+\frac{1}{x_{n}}$，$n=1,2,\cdots$，求$\lim_{n\to\infty}x_{n}$.
          \vspace{12em}

    \item 设数列$\{x_{n}\}_{n=1}^{\infty}$收敛，数列$\{y_{n}\}_{n=1}^{\infty}$发散，讨论$\{x_{n}+y_{n}\}_{n=1}^{\infty}$和$\left\{
              \frac{x_{n}}{y_{n}}\right\}_{n=1}^{\infty}$是否发散，若是，请证明之，若不是，请举反例.
          \vspace{10em}
\end{enumerate}

\section{证明题(每小题 5 分，共 10 分)}
\begin{enumerate}
    \setcounter{enumi}{15}
    \item 设函数$f(x)$在$[a,b]$上连续，且$f(a)=f(b)=0$，$f'_{+}(a)f'_{-}(b)>0$，则在$(a,b)$内至少存在一点$\xi$使得$f(\xi)=0$.
          \vspace{10em}

    \item 证明不等式
          \[
              \frac{1}{\sqrt{\left(1+x^{2}\right)\left(1+y^{2}\right)}}<\frac{\arctan y-\arctan x}{y-x}<1\qquad(0<x<y).
          \]
          \vspace{10em}
\end{enumerate}

\chapter{2026-2027学年微积分（B）（上）期中考试参考答案}
\section{基本计算题(每小题 6 分，共 60 分)}
\begin{enumerate}
    \item \textbf{Solution}. 记$A_{n}=\frac{1!+2!+3!+\cdots+n!}{n!}$（$n\geq3$）.显然$A_{n}\geq1$，且
          \[
              \begin{aligned}
                  A_{n}= {} & 1+\frac{(n-1)!}{n!}+\frac{(n-2)!+\cdots+2!+1!}{n!}                     \\
                  \leq {}   & 1+\frac{1}{n}+\frac{(n-2)(n-2)!}{n!}=1+\frac{1}{n}+\frac{n-2}{n(n-1)},
              \end{aligned}
          \]
          而$\lim_{n\to\infty}\left(1+\frac{1}{n}+\frac{n-2}{n(n-1)}\right)=1$，由夹逼准则得
          \[
              l=\lim_{n\to\infty}A_{n}=1.
          \]

    \item \textbf{Solution}. 令$t=x-2$，则$x\to2$时$t\to0$，且$x-1=1+t$.

          当$t\to0$时，$(1+t)^{\frac{1}{k}}-1\sim\frac{t}{k}$，故
          \[
              1-\sqrt[k]{x-1}=1-(1+t)^{\frac{1}{k}}\sim-\frac{t}{k},\qquad k=2,3,\cdots,n.
          \]
          因此
          \[
              (1-\sqrt{x-1})(1-\sqrt[3]{x-1})\cdots(1-\sqrt[n]{x-1})
              \sim\left(-\frac{t}{2}\right)\left(-\frac{t}{3}\right)\cdots\left(-\frac{t}{n}\right)
              =\frac{(-1)^{n-1}}{n!}t^{n-1},
          \]
          于是
          \[
              l=\lim_{t\to0}\frac{\frac{(-1)^{n-1}}{n!}t^{n-1}}{t^{n-1}}=\frac{(-1)^{n-1}}{n!}.
          \]

    \item \textbf{Solution}. 当$x\to0$时，$\tan^{2}x\sim x^{2}$，$1-\cos x^{2}\sim\frac{x^{4}}{2}$，且
          \[
              1-\sqrt{\cos x}=\frac{1-\cos x}{1+\sqrt{\cos x}}\sim\frac{x^{2}}{4}.
          \]
          所以
          \[
              l=\lim_{x\to0}\frac{x^{2}}{\frac{x^{4}}{2}}\cdot\frac{x^{2}}{4}=\frac{1}{2}.
          \]

    \item \textbf{Solution}. 因为
          \[
              \frac{x^{2}+1}{x+1}=\frac{(x+1)(x-1)+2}{x+1}=x-1+\frac{2}{x+1},
          \]
          所以
          \[
              \lim_{x\to\infty}\left(\frac{x^{2}+1}{x+1}-ax-b\right)=\lim_{x\to\infty}\left[(1-a)x-(1+b)+\frac{2}{x+1}\right]=\frac{1}{2},
          \]
          故必有$1-a=0$，$-(1+b)=\frac{1}{2}$，即$a=1$，$b=-\frac{3}{2}$.

    \item \textbf{Solution}. 令$t=\tan x$，则$x\to0$时$t\to0$，且$t\sim x$.因为
          \[
              \begin{aligned}
                  \sin(\tan x)-\tan(\tan x)= {} & \sin t-\tan t=\tan t\left(\cos t-1\right)                                 \\
                  \sim {}                       & t\cdot\left(-\frac{t^{2}}{2}\right)=-\frac{t^{3}}{2}\sim-\frac{x^{3}}{2},
              \end{aligned}
          \]
          所以无穷小量$\sin(\tan x)-\tan(\tan x)$的主部为$-\frac{1}{2}x^{3}$，阶数为3.

    \item \textbf{Solution}. 在$x=1$处，$f(1)=\mathrm{e}$，
          \[
              \lim_{x\to1^{-}}f(x)=\lim_{x\to1^{-}}\mathrm{e}^{x}=\mathrm{e},\qquad
              \lim_{x\to1^{+}}f(x)=\lim_{x\to1^{+}}\frac{2}{\pi}\arctan\frac{1}{x-1}=\frac{2}{\pi}\cdot\frac{\pi}{2}=1.
          \]
          因为$f(1^{-})\neq f(1^{+})$，故$x=1$是$f(x)$的跳跃间断点，$f(x)$在$x=1$处不连续，从而在$x=1$处不可导.

          当$x<1$时，$f'(x)=\mathrm{e}^{x}$；当$x>1$时，
          \[
              f'(x)=\frac{2}{\pi}\cdot\frac{1}{1+\frac{1}{(x-1)^{2}}}\cdot\left(-\frac{1}{(x-1)^{2}}\right)
              =-\frac{2}{\pi\left[1+(x-1)^{2}\right]}.
          \]
          所以$f(x)$在$x=1$处不可导，在其余点处可导，且
          \[
              f'(x)=
              \begin{cases}
                  \mathrm{e}^{x},                         & x<1, \\
                  -\frac{2}{\pi\left[1+(x-1)^{2}\right]}, & x>1.
              \end{cases}
          \]

    \item \textbf{Solution}. 由极坐标与直角坐标的关系，
          \[
              \begin{cases}
                  x=r\cos\theta=\mathrm{e}^{\theta}\cos\theta, \\
                  y=r\sin\theta=\mathrm{e}^{\theta}\sin\theta,
              \end{cases}
          \]
          求导得
          \[
              x'_{\theta}=\mathrm{e}^{\theta}\left(\cos\theta-\sin\theta\right),\quad
              y'_{\theta}=\mathrm{e}^{\theta}\left(\sin\theta+\cos\theta\right),\quad
              x''_{\theta}=-2\mathrm{e}^{\theta}\sin\theta,\quad
              y''_{\theta}=2\mathrm{e}^{\theta}\cos\theta.
          \]
          由参数方程所确定的函数的二阶导数公式，
          \[
              \frac{\mathrm{d}^{2}y}{\mathrm{d}x^{2}}=\frac{x'_{\theta}y''_{\theta}-x''_{\theta}y'_{\theta}}{\left(x'_{\theta}\right)^{3}}
              =\frac{2\mathrm{e}^{2\theta}\left[\left(\cos\theta-\sin\theta\right)\cos\theta+\sin\theta\left(\sin\theta+\cos\theta\right)\right]}{\mathrm{e}^{3\theta}\left(\cos\theta-\sin\theta\right)^{3}}
              =\frac{2\mathrm{e}^{-\theta}}{\left(\cos\theta-\sin\theta\right)^{3}}.
          \]

    \item \textbf{Solution}. 当$x\neq1$时，
          \[
              \frac{x^{2}-1}{x^{2}-3x+2}=\frac{(x-1)(x+1)}{(x-1)(x-2)}=\frac{x+1}{x-2}.
          \]
          $f(x)$在$x=0,1,2$处无定义，故$x=0,1,2$都是$f(x)$的间断点.

          在$x=1$处，$\lim_{x\to1}f(x)=\lim_{x\to1}\frac{x+1}{x-2}+\sin1=-2+\sin1$，极限存在，所以$x=1$是$f(x)$的可去间断点；

          在$x=2$处，$\lim_{x\to2^{-}}f(x)=-\infty$，$\lim_{x\to2^{+}}f(x)=+\infty$，所以$x=2$是$f(x)$的无穷间断点（第二类间断点）；

          在$x=0$处，取$x_{n}=\frac{1}{2n\pi}$，$x'_{n}=\frac{1}{2n\pi+\frac{\pi}{2}}$，则$x_{n},x'_{n}\to0\ (n\to\infty)$，而
          \[
              \lim_{n\to\infty}f(x_{n})=-\frac{1}{2},\qquad\lim_{n\to\infty}f(x'_{n})=\frac{1}{2},
          \]
          所以$\lim_{x\to0}f(x)$不存在，且$f(x)$在$x=0$的某去心邻域内有界，

          故$x=0$是$f(x)$的振荡间断点，是第二类间断点.

    \item \textbf{Solution}. （1）当$x\neq0$时，$\left|f(x)\right|=\left|x^{\alpha}\sin\frac{1}{x}\right|\leq|x|^{\alpha}$.

          当$\alpha>0$时，$\lim_{x\to0}|x|^{\alpha}=0$，由夹逼准则得$\lim_{x\to0}f(x)=0=f(0)$，即$f(x)$在$x=0$处连续；

          当$\alpha=0$时，$f(x)=\sin\frac{1}{x}\ (x\neq0)$，$\lim_{x\to0}f(x)$不存在；当$\alpha<0$时，取$x_{n}=\frac{1}{2n\pi+\frac{\pi}{2}}$，

          则$x_{n}\to0$，而$f(x_{n})=x_{n}^{\alpha}\to+\infty$，$\lim_{x\to0}f(x)$不存在.

          因此，当且仅当$\alpha>0$时，$f(x)$在$x=0$处连续.

          （2）$f'(0)=\lim_{x\to0}\frac{f(x)-f(0)}{x-0}=\lim_{x\to0}x^{\alpha-1}\sin\frac{1}{x}$.

          当$\alpha>1$时，$\left|x^{\alpha-1}\sin\frac{1}{x}\right|\leq|x|^{\alpha-1}\to0$，故$f'(0)=0$，即$f(x)$在$x=0$处可导；

          当$\alpha=1$时，$x^{\alpha-1}\sin\frac{1}{x}=\sin\frac{1}{x}$，极限不存在；当$\alpha<1$时，取$x_{n}=\frac{1}{2n\pi+\frac{\pi}{2}}$，

          则$x_{n}^{\alpha-1}\to+\infty$，极限不存在.

          因此，当且仅当$\alpha>1$时，$f(x)$在$x=0$处可导.

          \textbf{注}：当$x<0$时$x^{\alpha}$未必有定义，以上讨论默认其有定义，这不影响上述结论.

    \item \textbf{Solution}. 因为$f(x)$是正值函数，故
          \[
              y=\left[f\left(x^{2}\right)\right]^{\frac{1}{x}}=\mathrm{e}^{\frac{\ln f\left(x^{2}\right)}{x}}.
          \]
          两边取对数得$\ln y=\dfrac{\ln f\left(x^{2}\right)}{x}$，两边关于$x$求导，
          \[
              \frac{y'}{y}=\frac{\frac{2xf'\left(x^{2}\right)}{f\left(x^{2}\right)}\cdot x-\ln f\left(x^{2}\right)}{x^{2}},
          \]
          所以
          \[
              y'=y\cdot\frac{2x^{2}f'\left(x^{2}\right)-f\left(x^{2}\right)\ln f\left(x^{2}\right)}{x^{2}f\left(x^{2}\right)}.
          \]
          当$x=1$时，$y=f(1)=2$，于是
          \[
              \left.y'\right|_{x=1}=2\cdot\frac{2f'(1)-f(1)\ln f(1)}{f(1)}=2\cdot\frac{2-2\ln2}{2}=2\left(1-\ln2\right),
          \]
          故
          \[
              \left.\mathrm{d}y\right|_{x=1}=2\left(1-\ln2\right)\mathrm{d}x.
          \]
\end{enumerate}

\section{综合题(每小题 6 分，共 30 分)}
\begin{enumerate}
    \setcounter{enumi}{10}
    \item \textbf{Solution}. 所给条件即为
          \[
              \lim_{x\to0}\frac{x\sin x+f(x)}{x^{2}}=3,
          \]
          故
          \[
              f(x)=3x^{2}-x\sin x+o\left(x^{2}\right)=3x^{2}-x^{2}+o\left(x^{2}\right)=2x^{2}+o\left(x^{2}\right)\qquad(x\to0).
          \]
          因此$\lim_{x\to0}f(x)=0$，由$f(x)$在$x=0$处连续得$f(0)=0$；
          \[
              f'(0)=\lim_{x\to0}\frac{f(x)-f(0)}{x-0}=\lim_{x\to0}\frac{2x^{2}+o\left(x^{2}\right)}{x}=0,
          \]
          且
          \[
              \lim_{x\to0}\frac{f(x)}{x^{2}}=\lim_{x\to0}\frac{2x^{2}+o\left(x^{2}\right)}{x^{2}}=2.
          \]

    \item \textbf{Solution}. 因$f(0)=1$，故$y=1$对应$x=0$，即$g(1)=0$.

          对方程$f^{3}(x)+f(x)=x+2$两边关于$x$求导，得
          \[
              3f^{2}(x)f'(x)+f'(x)=1,\qquad\text{即}\quad f'(x)\left[3f^{2}(x)+1\right]=1,
          \]
          所以
          \[
              f'(x)=\frac{1}{3f^{2}(x)+1},\qquad f'(0)=\frac{1}{3\times1^{2}+1}=\frac{1}{4}.
          \]
          把$f'(x)\left[3f^{2}(x)+1\right]=1$两边再关于$x$求导，得
          \[
              f''(x)\left[3f^{2}(x)+1\right]+6f(x)\left[f'(x)\right]^{2}=0,
          \]
          所以
          \[
              f''(0)=-\frac{6f(0)\left[f'(0)\right]^{2}}{3f^{2}(0)+1}=-\frac{6\times\frac{1}{16}}{4}=-\frac{3}{32}.
          \]
          因$f'(x)>0$，反函数$x=g(y)$可导，由反函数的求导法则及$\frac{\mathrm{d}x}{\mathrm{d}y}=\frac{1}{f'(x)}$，
          \[
              \frac{\mathrm{d}^{2}x}{\mathrm{d}y^{2}}=\frac{\mathrm{d}}{\mathrm{d}y}\left(\frac{1}{f'(x)}\right)
              =-\frac{f''(x)}{\left[f'(x)\right]^{2}}\cdot\frac{\mathrm{d}x}{\mathrm{d}y}
              =-\frac{f''(x)}{\left[f'(x)\right]^{3}}.
          \]
          故
          \[
              \left.\frac{\mathrm{d}^{2}x}{\mathrm{d}y^{2}}\right|_{y=1}=-\frac{f''(0)}{\left[f'(0)\right]^{3}}
              =-\frac{-\frac{3}{32}}{\frac{1}{64}}=6.
          \]

    \item \textbf{Solution}. 设时刻$t$漏斗中水的深度为$h$，桶中水的深度为$H$，漏斗内水面的半径为$r$，则$r=\dfrac{6}{24}h=\dfrac{h}{4}$.

          漏斗中水的体积与桶中水的体积分别为
          \[
              V_{1}=\frac{1}{3}\pi r^{2}h=\frac{\pi h^{3}}{48},\qquad V_{2}=\pi\cdot6^{2}\cdot H=36\pi H.
          \]
          两式两边关于$t$求导，得
          \[
              \frac{\mathrm{d}V_{1}}{\mathrm{d}t}=\frac{\pi h^{2}}{16}\frac{\mathrm{d}h}{\mathrm{d}t},\qquad
              \frac{\mathrm{d}V_{2}}{\mathrm{d}t}=36\pi\frac{\mathrm{d}H}{\mathrm{d}t}.
          \]
          因为漏斗中减少的水全部流入桶内，所以$\dfrac{\mathrm{d}V_{2}}{\mathrm{d}t}=-\dfrac{\mathrm{d}V_{1}}{\mathrm{d}t}$.

          当$h=16\mathrm{cm}$，$\dfrac{\mathrm{d}h}{\mathrm{d}t}=-2\mathrm{cm/s}$时，
          \[
              \frac{\mathrm{d}V_{1}}{\mathrm{d}t}=\frac{\pi\cdot16^{2}}{16}\times(-2)=-32\pi,
          \]
          于是
          \[
              \frac{\mathrm{d}H}{\mathrm{d}t}=\frac{1}{36\pi}\cdot32\pi=\frac{8}{9}.
          \]
          即此时桶中水面上升的速度为$\dfrac{8}{9}\mathrm{cm/s}$.

    \item \textbf{Solution}. 由$x_{1}=1$及$x_{n+1}=1+\dfrac{1}{x_{n}}$，用数学归纳法易知$x_{n}\geq1\ (n=1,2,\cdots)$.

          设$l=\dfrac{1+\sqrt{5}}{2}$，则$l>1$，且$l^{2}-l-1=0$，即$l=1+\dfrac{1}{l}$.于是
          \[
              x_{n+1}-l=\frac{1}{x_{n}}-\frac{1}{l}=\frac{l-x_{n}}{x_{n}l},
          \]
          故
          \[
              \left|x_{n}-l\right|=\frac{1}{x_{n}l}\left|x_{n-1}-l\right|\leq\frac{1}{l}\left|x_{n-1}-l\right|
              \leq\cdots\leq\frac{1}{l^{n-1}}\left|x_{1}-l\right|.
          \]
          因为$l>1$，$\lim_{n\to\infty}\frac{\left|x_{1}-l\right|}{l^{n-1}}=0$，由夹逼准则得$\lim_{n\to\infty}\left|x_{n}-l\right|=0$，所以
          \[
              \lim_{n\to\infty}x_{n}=l=\frac{1+\sqrt{5}}{2}.
          \]

    \item \textbf{Solution}. 设$\lim_{n\to\infty}x_{n}=a$.

          （1）$\{x_{n}+y_{n}\}$必发散.用反证法：若$\{x_{n}+y_{n}\}$收敛，设$\lim_{n\to\infty}\left(x_{n}+y_{n}\right)=b$，则
          \[
              \lim_{n\to\infty}y_{n}=\lim_{n\to\infty}\left[\left(x_{n}+y_{n}\right)-x_{n}\right]=b-a,
          \]
          这与$\{y_{n}\}$发散矛盾，所以$\{x_{n}+y_{n}\}$发散.

          （2）$\left\{\dfrac{x_{n}}{y_{n}}\right\}$不一定发散.例如取$x_{n}\equiv1$，$y_{n}=n$，则$\{x_{n}\}$收敛，$\{y_{n}\}$发散，而
          \[
              \lim_{n\to\infty}\frac{x_{n}}{y_{n}}=\lim_{n\to\infty}\frac{1}{n}=0,
          \]
          即$\left\{\dfrac{x_{n}}{y_{n}}\right\}$收敛.
\end{enumerate}

\section{证明题(每小题 5 分，共 10 分)}
\begin{enumerate}
    \setcounter{enumi}{15}
    \item \textbf{Proof}. 不妨设$f'_{+}(a)>0$，$f'_{-}(b)>0$.若两者均为负，则用$-f(x)$代替$f(x)$，下述论证同样适用.

          因为
          \[
              f'_{+}(a)=\lim_{x\to a^{+}}\frac{f(x)-f(a)}{x-a}=\lim_{x\to a^{+}}\frac{f(x)}{x-a}>0,
          \]
          由极限的局部保号性，存在$\delta_{1}>0$（不妨设$\delta_{1}<\dfrac{b-a}{2}$），使得当$x\in\left(a,a+\delta_{1}\right)$时$\dfrac{f(x)}{x-a}>0$，

          从而$f(x)>0$. 取定$x_{1}\in\left(a,a+\delta_{1}\right)$，则$f(x_{1})>0$，且$x_{1}<\dfrac{a+b}{2}$.

          同理，因为
          \[
              f'_{-}(b)=\lim_{x\to b^{-}}\frac{f(x)-f(b)}{x-b}=\lim_{x\to b^{-}}\frac{f(x)}{x-b}>0,
          \]
          存在$\delta_{2}>0$（不妨设$\delta_{2}<\dfrac{b-a}{2}$），使得当$x\in\left(b-\delta_{2},b\right)$时$\dfrac{f(x)}{x-b}>0$，此时$x-b<0$，从而$f(x)<0$.

          取定$x_{2}\in\left(b-\delta_{2},b\right)$，则$f(x_{2})<0$，且$x_{2}>\dfrac{a+b}{2}$.

          于是$a<x_{1}<x_{2}<b$，$f(x_{1})f(x_{2})<0$，且$f(x)$在$\left[x_{1},x_{2}\right]$上连续，

          由零点定理，至少存在一点$\xi\in\left(x_{1},x_{2}\right)\subset(a,b)$，使得
          \[
              f(\xi)=0.
          \]

    \item \textbf{Proof}. 设$f(t)=\arctan t$，则$f(t)$在$[x,y]$上连续，在$(x,y)$内可导，

          由Lagrange中值定理，至少存在一点$\xi\in(x,y)$，使得
          \[
              \frac{\arctan y-\arctan x}{y-x}=f'(\xi)=\frac{1}{1+\xi^{2}}.
          \]
          因为$\xi>x>0$，所以$\dfrac{1}{1+\xi^{2}}<1$，右边的不等式成立.

          另一方面，因为$0<x<\xi<y$，所以$1+\xi^{2}<1+y^{2}$，又$\sqrt{1+x^{2}}>1$，故
          \[
              1+\xi^{2}<1+y^{2}<\left(1+y^{2}\right)\sqrt{1+x^{2}}=\sqrt{\left(1+x^{2}\right)\left(1+y^{2}\right)},
          \]
          从而
          \[
              \frac{1}{1+\xi^{2}}>\frac{1}{\sqrt{\left(1+x^{2}\right)\left(1+y^{2}\right)}}.
          \]
          综上所述，
          \[
              \frac{1}{\sqrt{\left(1+x^{2}\right)\left(1+y^{2}\right)}}<\frac{\arctan y-\arctan x}{y-x}<1\qquad(0<x<y).
          \]
\end{enumerate}
\end{document}
